\documentclass[../../main.tex]{subfiles}
\begin{document}

{\bf [解]}

\paragraph{(1)}

$z=x+iy$（$x,y\in\mathbb{R}$，$y>0$）とおいて$f(z)$を計算すると
\begin{align}
  \begin{split}
    f(z)&=1-\dfrac{q}{z+1}\\
    &=1-\dfrac{q(\overline{z}+1)}{|z+1|^2}\\
    &=1-\dfrac{q(x+1)}{|z+1|^2}+\dfrac{qy}{|z+1|^2}i．
  \end{split}
\end{align}
である．ここで虚部について$|z+1|^2=(x+1)^2+y^2>0$および$qy>0$だから，$f(z)$の虚部は正である．従って，$f(z)\in H$である．

\paragraph{(2)}

$w=f(z)=\dfrac{z+1-q}{z+1}$とおく．$q>0$から$w\ne1$であり，
\begin{align}
  z=-1+\dfrac{q}{1-w}
\end{align}
となる．よって$f$の逆関数は
\begin{align}
  f^{-1}(z)=-1+\dfrac{q}{1-z}
  =\dfrac{-z+1-q}{z-1}
  \label{eq:1}
\end{align}
である．

(1)と同様に，$z=x+iy\in H$に対し
\begin{align}
  \operatorname{Im} f^{-1}(z)=\dfrac{qy}{|1-z|^2}>0
\end{align}
だから，$f^{-1}(z)\in H$である．

$n\ge2$のとき，$f_n(z)$に$f^{-1}$を作用させると
\begin{align}
  f^{-1}(f_n(z))=f^{-1}(f(f_{n-1}(z)))=f_{n-1}(z) \label{eq:2}
\end{align}
であり，また，$f_1(z)=f(z)$に$f^{-1}$を作用させると
\begin{align}
  f^{-1}(f_1(z))=z \label{eq:3}
\end{align}
である．

従って，$f_{10}(z)=f_5(z)$に$f^{-1}$を5回施して
\begin{align}
  f_{10}(z)=f_5(z)\quad(\forall z\in H)
  \quad\Longleftrightarrow\quad
  f_5(z)=z\quad(\forall z\in H)
\end{align}
となり，さらに1回施すと求める条件は
\begin{align}
  f_4(z)=f^{-1}(z)\quad(\forall z\in H) \label{eq:4}
\end{align}
と同値である．

さて，$f_4(z)$を直接計算して
\begin{align}
  f_2(z)&=\dfrac{(2-q)z+2(1-q)}{2z+2-q}，\\
  f_4(z)=f_2(f_2(z))
  &=\dfrac{(q^2-8q+8)z+4(1-q)(2-q)}{4(2-q)z+q^2-8q+8} \label{eq:5}
\end{align}
である．\cref{eq:4}において\cref{eq:1,eq:5}を比較してこれらが任意の$z$で一致する条件を考える．
\begin{align}
  \dfrac{-z+1-q}{z-1} = \dfrac{(q^2-8q+8)z+4(1-q)(2-q)}{4(2-q)z+q^2-8q+8}
\end{align}
で分母を払って
\begin{align}
  \left(4(2-q)z+q^2-8q+8\right)\left(-z+1-q\right) = \left((q^2-8q+8)z+4(1-q)(2-q)\right)\left(z-1\right)
\end{align}
が恒等式になれば良い．
ここで簡単のため$A=q^2-8q+8$，$B=4(2-q)$とおけば，
\begin{align}
  & \left(Bz+A\right)\left(-z+1-q\right) = \left(Az+B(1-q)\right)\left(z-1\right) \\
  & \left(Bz+A+Az+B(1-q)\right)(1-z) - q\left(Bz+A\right) = 0 \\
  & \left[(A+B)(z+1) -Bq\right](1-z) - q\left(Bz+A\right) = 0 \\
  & (A+B)(1-z^2) -(A+B)q = 0 \\
  & (A+B)(1-q-z^2) = 0
\end{align}
である．これが恒等式になる必要十分条件は
\begin{align}
  A+B=0
  \quad\Longleftrightarrow\quad
  q^2-8q+8=-4(2-q)
  \quad\Longleftrightarrow\quad
  q^2-12q+16=0
\end{align}
である．これを解いて
\begin{align}
  q=6\pm2\sqrt{5}
\end{align}
を得る．いずれも$q>0$を満たすからこれが答えである．

\end{document}
